% Appendix B --- computational notes.
% Sources: 02_replicator_suppressor.tex:542-677, trimmed from 135 lines to
% the mathematical content (b_hat, a_hat, the assembly of A and c, and the
% solve); the array-plumbing functions concatA, concatC, full_A, full_C and
% get2D_dist are deleted.  Quadtree note from 04:1-3 and :8-28, demoted from
% a numbered chapter section.
% The listing's variable names are the chapter's: the source used `alpha`
% for the birth rate and `beta` for the suppression rate, which invert the
% chapter's letters.

\section{Computational notes}
\label{path:app:computational}

\subsection{The first-step solver}
\label{path:app:solver}

The system \cref{path:eq:matrixsystem} is assembled directly from the jump
probabilities of \cref{path:eq:firststep} and solved once. The unknowns are
$\pi_{i,j}$ for $1\le i\le j\le N$; the two boundary conditions
$\pi_{0,j}=1$ and $\pi_{i,j}=0$ for $i>j$ enter as the inhomogeneous term
and as an omission from the index set respectively. The assembly is direct
enough to describe in words, and the working implementation is a file rather
than a printed page.

Index the transient states $(i,j)$ with $1\le i\le j\le N$ in any fixed
order. For each such state, place $\hat\alpha_j$ in the column of
$\pi_{i+1,j}$ whenever $i+1\le j$, and $\hat\beta_j$ in the column of
$\pi_{i-1,j-1}$ whenever $i-1\ge1$; when $i=1$ the second coefficient has no
column to go in, because $\pi_{0,j-1}=1$ is known, and it becomes that row's
entry of $\mathbf{c}$. Nothing else is non-zero. One call to a dense
linear solver on $(\mathbf{I}-\mathbf{A})\mathbf{p}^{(0)}=\mathbf{c}$
then returns every $\pi_{i,j}$ at once. At $N=20$ the system has $210$
unknowns and solves in milliseconds, so the block-triangular structure of
\cref{path:fig:levels} is a fact about the matrix rather than a
computational necessity; it becomes one only for much larger stores.

Two features of that recipe are worth noting, and both come from
\cref{path:sec:rsfirststep}. The coefficients depend on $j$ alone, so each
takes one argument where the state has two; and the omission of the states
$i>j$ is not an approximation but \cref{path:prop:certainty}, which is what
keeps the system finite despite $X_t$ being unbounded. Get either wrong and
the solve still returns numbers, which is why the implementation checks its
diagonal against the closed form of \cref{path:prop:diagonal} on every run,
and agrees to $1.9\times10^{-16}$.

The working implementation sits with the chapter's figure-generation
scripts. It is self-contained, depends only on \texttt{numpy} and
\texttt{matplotlib}, carries the assembly above, the solve, that check, and
the plotting that produces \cref{path:fig:repsupp20,path:fig:repsupp10}.

\subsection{Quadtree sampling}
\label{path:app:quadtree}

Collision tests for the particle runs of
\cref{path:fig:basic,path:fig:rs,path:fig:loads} were
quadtree-indexed, so that only agents local to a macrophage are tested.
The index is rebuilt each frame; the saving is that collision testing stops
being quadratic in the particle count, which is what the large-load runs of
\cref{path:fig:loads} require.

\begin{figure}[htbp]
    \centering
    \begin{subfigure}{0.40\textwidth}
        \centering
        \includegraphics[width=\textwidth]{QT1.png}
        \caption{}
    \end{subfigure}
    \hspace{0.03\textwidth}
    \begin{subfigure}{0.40\textwidth}
        \centering
        \includegraphics[width=\textwidth]{QT2.png}
        \caption{}
    \end{subfigure}

    \caption{The quadtree index during a simulation run, with the cells
    local to the macrophage shaded pink; only the particles inside those
    cells are tested for collision. (a) A sparse field, where the index buys
    little. (b) A dense one, where it buys a great deal: the shaded region
    is a small fraction of the frame, so the number of collision tests grows
    far more slowly than the number of particles. This is why the
    large-load runs of \cref{path:fig:loads} are feasible at all.}
    \label{path:fig:quadtree}
\end{figure}

\Cref{path:fig:quadtree} shows the index in a sparse field and a dense one.

\FloatBarrier
